\chapter{Brainteasers and Logic}\label{ch:iv:brainteasers-and-logic}

A candidate is asked how a hundred people standing in a line, each seeing the hats of everyone in front, can
agree on a rule so that all but one name the colour of their own hat. He says ``this is the parity puzzle'' and
recites the rule for two colours. The interviewer changes it to three colours and waits. The puzzle was never
the point: the family was, and the idea behind the family, and a candidate who knows only the answer to the
famous version has nothing to say when the version changes. This chapter organises brainteasers into families,
gives the idea behind each, and then asks only variants.

\section{Invariants and parity}

\begin{definition}[Invariant method]\label{def:iv:brainteasers-and-logic:invariant}
The \emph{invariant method}\index{invariant method} solves a puzzle about a process by finding a quantity that
no allowed move changes (a parity, a sum modulo $m$, a product, a colouring count); every reachable state
shares the starting value, so a state with another value is unreachable, and the final state is determined.
\end{definition}

\begin{method}[Finding an invariant]\label{met:iv:brainteasers-and-logic:invariant}
\begin{enumerate}
\item Play a few moves by hand and write the state after each.
\item Try, in order: the parity of each count, the sum modulo small numbers, the sum or product of a function
of the entries ($\prod(1 + x_i)$ for merges of the form $a + b + ab$), a colouring of a board.
\item Prove that each move preserves it; then read the answer from the initial value.
\end{enumerate}
\end{method}

\begin{example}[The hat line]\label{ex:iv:brainteasers-and-logic:hats}
With $k$ colours numbered $0, \dots, k-1$, the last person in the line (who sees everyone else) announces the
sum of the colours in front modulo $k$. Each following person knows that sum, hears every answer given after
the first, sees everyone in front, and so can subtract to find their own colour. All but the first are right,
whatever $k$: the invariant is the announced sum.
\end{example}

\section{Working backwards}

\begin{definition}[Backward induction]\label{def:iv:brainteasers-and-logic:backward}
\emph{Backward induction}\index{backward induction} solves a finite sequential game or decision problem by
determining the best action at the last stage, then at the one before given the best continuation, and so on to
the first stage.
\end{definition}

In a take-away game (players alternately remove tokens, the one who takes the last wins), a position is losing
for the player to move if every move leads to a winning position for the other, and winning if some move leads
to a losing one. Computing this from zero upwards gives the pattern (\cref{fig:iv:brainteasers-and-logic:game}):
with moves of 1 or 2 the losing positions are the multiples of 3; with moves of 1 to $m$, the multiples of
$m+1$; with irregular move sets the pattern is eventually periodic but has to be computed.

\begin{omfigure}
\begin{tikzpicture}[font=\footnotesize,>={Stealth},
  w/.style={circle,draw=omDef,fill=omDef!6,minimum size=0.75cm},
  l/.style={circle,draw=omThm,fill=omThm!12,minimum size=0.75cm,thick}]
\foreach \n/\s in {7/w,6/l,5/w,4/w,3/l,2/w,1/w,0/l} {
  \node[\s] (p\n) at ({1.45*(7-\n)},0) {\n};
}
\foreach \a/\b in {7/6,6/5,5/4,4/3,3/2,2/1,1/0} { \draw[->,gray] (p\a) -- (p\b); }
\foreach \a/\b in {7/5,5/3,3/1} { \draw[->,gray] (p\a) to[bend left=35] (p\b); }
\foreach \a/\b in {6/4,4/2,2/0} { \draw[->,gray] (p\a) to[bend right=35] (p\b); }
\node[font=\scriptsize,anchor=west] at (-0.4,-1.25) {red: the player to move loses; blue: the player to move wins by moving to a red position};
\end{tikzpicture}

\omcaption{The take-away game with moves of one or two tokens, last to take wins, solved by \omterm{def:iv:brainteasers-and-logic:backward}{backward induction}
from 0. Arrows are moves (straight: take one; curved: take two). Positions 0, 3 and 6 are losing for the player
to move. Exhaustive search in \texttt{iv\_teasers.py}.}\label{fig:iv:brainteasers-and-logic:game}
\end{omfigure}

Division games (a senior proposer, a vote, a proposer who leaves if rejected) are solved the same way: find
the outcome with one player, then with two, and at each size give the minimum needed to the voters who would do
worst if the proposal failed. The classic version is Stewart's pirate puzzle (1999); the bank's variant has
different voting rules and a different answer.

\section{Strategy stealing, pigeonholes and extremal arguments}

\begin{definition}[Strategy-stealing argument]\label{def:iv:brainteasers-and-logic:stealing}
A \emph{strategy-stealing argument}\index{strategy-stealing argument} proves that the first player in a finite
game with no draws has a winning strategy, without finding it: if the second player had one, the first player
could make an arbitrary move and then follow that strategy as if second, and an extra move never hurts, which
is a contradiction.
\end{definition}

The argument proves existence, not construction: Gale's chocolate-bar game (1974), in which players remove a
rectangle from the corner of a bar and the player forced to take the poisoned square loses, is a first-player
win on every rectangle larger than one square, and nobody knows a general winning move. Gale's proof that Hex
cannot end in a draw (1979) is the other famous use.

Pigeonhole arguments show that something must coincide because there are more objects than boxes; extremal
arguments look at the largest or smallest element and show it has the property asked for. Both reward naming
the boxes or the extreme element explicitly.

\begin{example}[A block of trades summing to a multiple of ten]\label{ex:iv:brainteasers-and-logic:pigeonhole}
Among any ten integer P\&Ls in sequence, some block of consecutive ones sums to a multiple of ten. The eleven
prefix sums (including the empty one) fall into ten residues modulo ten, so two coincide, and the block between
them sums to a multiple of ten.
\end{example}

\section{Information puzzles}

Weighing and questioning puzzles are about counting outcomes: $k$ weighings on a balance have $3^k$ outcomes and
can distinguish at most $3^k$ cases, and a design that splits the cases as evenly as possible at each step
reaches the bound or nearly. Common-knowledge puzzles are about what each person knows about what the others
know: an announcement that everyone already knew can still change behaviour, because after it everyone knows that
everyone knows it (Fagin, Halpern, Moses and Vardi, 1995).

\begin{method}[Information puzzles]\label{met:iv:brainteasers-and-logic:info}
\begin{enumerate}
\item Count the cases to distinguish and the outcomes each test has; the logarithm is a lower bound on the
number of tests.
\item Design each test to split the remaining cases as evenly as the outcomes allow.
\item For knowledge puzzles, start from the smallest case (one person, one round) and induct: ``if there were
only one, she would know at once; nobody spoke, so there are at least two''.
\end{enumerate}
\end{method}

The simultaneous hat game, in which players guess or pass at the same time and win if at least one guesses and
nobody is wrong, was popularised by Ebert's thesis (1998); its solution for seven players uses a Hamming code.

\section{Worked answers}

\begin{example}[Pigeonholes, and the fewest that force it]\label{ex:iv:brainteasers-and-logic:gaps}
\emph{``Thirteen trades print within one hour, from 10:00 to 11:00 inclusive. Prove that two of them are at most five
minutes apart. Is thirteen the fewest that forces it?''} Cut the hour into twelve closed intervals of five minutes; with
thirteen trades, two share an interval and are at most five minutes apart. For the second part, an extremal
construction: twelve trades at 10:00, 10:05 and a hair, 10:10 and two hairs, and so on, span $11 \times (5 + \epsilon)$
minutes, which fits in the hour for small $\epsilon$, with every gap above five minutes. So twelve do not force it and
thirteen is the fewest. The two halves (the argument and the construction showing it is tight) are both expected;
the second is what separates a proof from a recitation.
\end{example}

\begin{example}[A parity invariant]\label{ex:iv:brainteasers-and-logic:diff}
\emph{``The numbers 1 to 20 are on a board. Repeatedly erase two numbers and write the absolute value of their
difference. Can the last number be 1? Can it be 0?''} Look for what an operation cannot change. Since $|a - b| \equiv a +
b \pmod 2$, the parity of the sum of the board never changes; it starts at $210$, even, so the last number is even and
cannot be 1. Zero is reachable: pair the numbers as $(1,2), (3,4), \dots, (19,20)$ to get ten 1s, then pair the 1s to
get five 0s, and every further difference of zeros is zero. The invariant rules outcomes out; an explicit sequence rules
one in, and the answer needs both.
\end{example}

\section{Question bank}

\begin{interviewq}[$\star$ \iqroles{trader} \iqfirm{market maker}]\label{iq:iv:brainteasers-and-logic:1}
A bag holds 20 red and 13 blue tokens. Repeatedly take two out: if they have the same colour put back one red,
if they differ put back one blue. What colour is the last token?
\end{interviewq}

\begin{interviewq}[$\star$ \iqroles{trader, developer} \iqfirm{proprietary firm}]\label{iq:iv:brainteasers-and-logic:2}
Thirty tokens; players alternately take one to four; whoever takes the last wins. Do you want to go first? With
31 tokens, what is your first move?
\end{interviewq}

\begin{interviewq}[$\star$ \iqroles{trader, researcher} \iqfirm{any}]\label{iq:iv:brainteasers-and-logic:3}
A firm employs 400 people. Prove that two of them have the same birthday.
\end{interviewq}

\begin{interviewq}[$\star$ \iqroles{researcher, developer} \iqfirm{any}]\label{iq:iv:brainteasers-and-logic:4}
In a group of traders, some pairs have traded with each other. Prove that the number of traders who have traded
with an odd number of the others is even.
\end{interviewq}

\begin{interviewq}[$\star$ \iqroles{trader} \iqfirm{market maker}]\label{iq:iv:brainteasers-and-logic:5}
You have a 3-minute and a 5-minute sand timer. Measure exactly 7 minutes.
\end{interviewq}

\begin{interviewq}[$\star\star$ \iqroles{trader, developer} \iqfirm{proprietary firm}]\label{iq:iv:brainteasers-and-logic:6}
A hundred people stand in a line, each seeing the hats of all those in front; hats come in three colours. Guessing
from the back, each says one colour aloud. Give a rule agreed in advance that guarantees at least 99 correct
answers, and explain why it works.
\end{interviewq}

\begin{interviewq}[$\star\star$ \iqroles{trader, researcher} \iqfirm{market maker}]\label{iq:iv:brainteasers-and-logic:7}
Ten bags of coins; in one bag every coin weighs 0.9 gram instead of 1 gram. With a single reading of a digital
scale, find the bag.
\end{interviewq}

\begin{interviewq}[$\star\star$ \iqroles{trader, researcher} \iqfirm{any}]\label{iq:iv:brainteasers-and-logic:8}
Four partners, ranked by seniority, divide 100 coins. The most senior remaining partner proposes a division;
it passes if at least half of the remaining partners (the proposer included) vote for it; if it fails, the
proposer leaves with nothing and the next proposes. Each partner votes for a proposal only if it gives strictly
more than they would get after a rejection. What does the most senior partner propose?
\end{interviewq}

\begin{interviewq}[$\star\star$ \iqroles{researcher, developer} \iqfirm{any}]\label{iq:iv:brainteasers-and-logic:9}
Two players alternately bite a rectangular corner off a four-by-five chocolate bar (choosing a square and
removing it together with every square below and to its right); whoever must take the top-left square loses.
Prove that the first player can win without finding the winning move.
\end{interviewq}

\begin{interviewq}[$\star\star$ \iqroles{trader, researcher} \iqfirm{proprietary firm}]\label{iq:iv:brainteasers-and-logic:10}
Five traders wear red or blue badges; each sees the others' badges but not their own; three badges are red. A
supervisor says: ``At least one of you wears red. Whoever knows their colour, step forward.'' She repeats the
instruction at each round. Who steps forward, and at which round? What did the supervisor's sentence add, since
everyone could already see a red badge?
\end{interviewq}

\begin{interviewq}[$\star\star\star$ \iqroles{researcher, developer} \iqfirm{systematic fund}]\label{iq:iv:brainteasers-and-logic:11}
Three players each get a red or blue hat with probability one half. Each sees the other two hats, then all
simultaneously guess their own colour or pass. The team wins if at least one guesses and nobody guesses wrong. Find
a strategy that wins three times in four and prove that no strategy does better.
\end{interviewq}

\begin{interviewq}[$\star\star\star$ \iqroles{trader, researcher} \iqfirm{any}]\label{iq:iv:brainteasers-and-logic:12}
The numbers 1 to 10 are written on a board. Repeatedly erase two numbers $a$ and $b$ and write $a + b + ab$.
What number is left at the end, and why does the order not matter?
\end{interviewq}

\begin{interviewq}[$\star\star\star$ \iqroles{trader, developer} \iqfirm{market maker}]\label{iq:iv:brainteasers-and-logic:13}
Twenty tokens; players alternately take one, three or four; whoever takes the last wins. Who wins, and what is the
first move? Describe the losing positions.
\end{interviewq}

\begin{interviewq}[$\star\star\star$ \iqroles{researcher} \iqfirm{systematic fund}]\label{iq:iv:brainteasers-and-logic:14}
Prove that among any ten integer daily P\&L figures, listed in order, some run of consecutive days has a total
that is a multiple of ten. Is ten days the fewest that guarantees it?
\end{interviewq}

\begin{omsources}
\item I.~Stewart, ``A puzzle for pirates'', \emph{Scientific American} 280(5), 1999, 98--99.
\item D.~Gale, ``A curious Nim-type game'', \emph{American Mathematical Monthly} 81(8), 1974, 876--879.
\item D.~Gale, ``The game of Hex and the Brouwer fixed-point theorem'', \emph{American Mathematical Monthly}
86(10), 1979, 818--827.
\item R.~Fagin, J.~Y.~Halpern, Y.~Moses and M.~Y.~Vardi, \emph{Reasoning about Knowledge}, MIT Press, 1995.
\item T.~Ebert, doctoral thesis, University of California, Santa Barbara, 1998 (the simultaneous hat game), as
reported in P.~Winkler, \emph{Mathematical Puzzles: A Connoisseur's Collection}, A K Peters, 2004.
\end{omsources}
